\clearpage
\VerdierPhysicalPageAnchor{0162}%
\begingroup\small\setlength{\parskip}{0pt}\setlength{\emergencystretch}{1.5em}
\VerdierSemanticAnchor{ch4}
\section*{4. スペクトル対象。}

[1] 第 XV 章 \S 7 におけるスペクトル系列の叙述を取り上げ、三角圏に
適合するよう改める。

\VerdierSemanticAnchor{ch4-4-1}
\subsection*{4.1. スペクトル対象の定義。}

\VerdierSemanticAnchor{ch4-4-1-1}
\medskip\noindent\textbf{4.1.1.}\quad
$\mathcal J$ を圏とし、各 $n\in\mathbb N$ に対して
$\mathcal{F}\ell_n(\mathcal J)$ を、$n$ 個の $\mathcal J$ の合成可能な
射の列からなる圏とする（第 I 章、3.3.6）。次の三つの関手
$s_i\colon\mathcal{F}\ell_2(\mathcal J)\to\mathcal{F}\ell_1(\mathcal J)$、
$1\leq i\leq3$ を用いる：
\VerdierSemanticAnchor{ch4-4-1-1-1}
\[
\left\{
\begin{aligned}
&\bigl(\mathord\bullet\xrightarrow f\mathord\bullet\xrightarrow g\mathord\bullet\bigr)
   \in\operatorname{Ob}\mathcal{F}\ell_2(\mathcal J),\\
&s_1\bigl(\mathord\bullet\xrightarrow f\mathord\bullet\xrightarrow g\mathord\bullet\bigr)
   =\bigl(\mathord\bullet\xrightarrow f\mathord\bullet\bigr),\\
&s_2\bigl(\mathord\bullet\xrightarrow f\mathord\bullet\xrightarrow g\mathord\bullet\bigr)
   =\bigl(\mathord\bullet\xrightarrow{gf}\mathord\bullet\bigr),\\
&s_3\bigl(\mathord\bullet\xrightarrow f\mathord\bullet\xrightarrow g\mathord\bullet\bigr)
   =\bigl(\mathord\bullet\xrightarrow g\mathord\bullet\bigr).
\end{aligned}
\right.
\tag{4.1.1.1}
\]
および二つの関手の射：
\VerdierSemanticAnchor{ch4-4-1-1-2}
\[
\left\{
\begin{aligned}
&u_1\colon s_1\longrightarrow s_2,
&&u_1(f,g)=
\begin{tikzcd}[ampersand replacement=\&,column sep=2.8em,row sep=1.9em,baseline=(current bounding box.center)]
{}\arrow[r,"f"]\arrow[d,"\operatorname{id}"']\&{}\arrow[d,"g"]\\
{}\arrow[r,"gf"']\&{}
\end{tikzcd},\\[2pt]
&u_2\colon s_2\longrightarrow s_3,
&&u_2(f,g)=
\begin{tikzcd}[ampersand replacement=\&,column sep=2.8em,row sep=1.9em,baseline=(current bounding box.center)]
{}\arrow[r,"gf"]\arrow[d,"f"']\&{}\arrow[d,"\operatorname{id}"]\\
{}\arrow[r,"g"']\&{}
\end{tikzcd}.
\end{aligned}
\right.
\tag{4.1.1.2}
\]

\VerdierSemanticAnchor{ch4-4-1-2}
\medskip\noindent\textbf{4.1.2.}\quad
$\mathcal D$ を三角圏、$\mathcal J$ を圏とする。
\emph{型 $\mathcal J$ で $\mathcal D$ に値をとるスペクトル対象}とは、
次の二つの関手からなる対象である：

a) 関手
$X\colon\mathcal{F}\ell_1(\mathcal J)\longrightarrow\mathcal D$：
\[
\bigl(\mathord\bullet\xrightarrow f\mathord\bullet\bigr)\longmapsto X_f.
\]
\endgroup
